\begin{ThreePartTable}

    \renewcommand\TPTminimum{\textwidth}

    \setlength\LTleft{0pt}
    \setlength\LTright{0pt}
    \setlength\tabcolsep{3pt}

    \begin{TableNotes}
        \item \textit{Nota.} En todas las operaciones, la posición inicial corresponde a la posición final de la operación anterior o, en ausencia de una operación previa, a la posición original del modelo tridimensional.
    \end{TableNotes}

    \begin{longtable}{@{}P{0.25\textwidth}P{0.18\textwidth}P{0.57\textwidth}@{}}

        \caption{Operaciones disponibles para la creación de animaciones de mantenimiento}
        \label{tab:operations} \\

        \toprule
        \textbf{Operación} &
        \textbf{Grados de libertad} &
        \textbf{Descripción} \\
        \toprule
        \endfirsthead

        \multicolumn{3}{r}{\textit{(Continuación de la Tabla \ref{tab:operations})}} \\
        \toprule
        \textbf{Operación} &
        \textbf{Grados de libertad} &
        \textbf{Descripción} \\
        \toprule
        \endhead

        \midrule[\heavyrulewidth]
        \multicolumn{3}{r}{\textit{Continúa en la siguiente página}}\\
        \endfoot

        \bottomrule
        \insertTableNotes\\
        \endlastfoot

        Punto a punto &
        6 GDL &
        Permite definir una posición y orientación finales para un submodelo tridimensional. \\

        \midrule

        Cilíndrica &
        2 GDL &
        Describe un desplazamiento a lo largo de una dirección determinada y una rotación final sobre ese mismo eje. \\

        \midrule

        Junta &
        3 GDL &
        Permite modificar únicamente la orientación del submodelo mediante tres grados de libertad rotacionales, manteniendo como referencia la posición final de la operación anterior. \\

        \midrule

        Tornillo &
        2 GDL &
        Representa un movimiento helicoidal definido por un eje de desplazamiento y un número de revoluciones alrededor de dicho eje. \\

    \end{longtable}

\end{ThreePartTable}